\section{Constructing new measure preserving systems from old ones}\label{sec:constr-new-measure-pres-sys-old-ones}

\subsection{Product maps}\label{sec:product-maps}

Let $(X_{1},\B_{1},\mu_{1})$ and $(X_{2},\B_{2},\mu_{2})$ be two $\sigma$-finite measure spaces. The \textbf{product $\sigma$-algebra} $\B_{1}\otimes\B_{2}$ on $X_{1}\times X_{2}$ is defined as the $\sigma$-algebra generated by the family of \emph{rectangles} $\{E_{1}\times E_{2} : E_{i}\in\B_{i},\ i\in\{1,2\}\}$. Then, one shows that there exists a unique measure $\mu_{1}\otimes\mu_{2}$ on the measurable space $(X_{1}\times X_{2},\B_{1}\otimes\B_{2})$, called the \emph{product measure}, such that
\begin{displaymath}
  \mu_{1}\otimes\mu_{2}(E_{1}\times E_{2})=\mu_{1}(E_{1})\mu_{2}(E_{2}), \quad\forall E_{1}\in\B_{1},\ \forall E_{2}\in\B_{2}.
\end{displaymath}

When $f_{i}\colon (X_{1},\B_{i},\mu_{i})\carr$, for $i\in\{1,2\}$, are measure preserving maps, then invoking Theorem~\ref{thm:uniqueness-measures} one can easily show that the measure $\mu_{1}\otimes\mu_{2}$ is $f_{1}\times f_{2}$-invariant, where $f_{1}\times f_{2}\colon X_{1}\times X_{2}\ni (x_{1},x_{2})\mapsto \big(f_{1}(x_{1}),f_{2}(x_{2})\big)\in X_{1}\times X_{2}$.

By induction, this construction can be easily extended from two to any (finite) number of factor maps. Moreover, this construction can be extended to countably many factors, assuming that each factor is a probability space (see~\cite[Appendix A.2]{VianaOliveiraFoundErgTh} for details).


\subsection{Induction or renormalization}\label{sec:induct-or-renormalization}

Let $(X,\B,\mu)$ be a probability space, $f\colon X\carr$ be a measure preserving map and $A\in\B$ a measurable set with $\mu(A)>0$. By Poincaré recurrence theorem (Corollary~\ref{cor:return-infinitely-many-times-Poincare}) we know that $\mu$-almost every point $x\in A$ returns to $A$ infinitely many times. So we can define the function $\rho_{A}\colon A\to\N$ by
\begin{displaymath}
  \rho_{A}(x):=\min\left\{n\in\N : f^{n}(x)\in A\right\},
\end{displaymath}
and this function is well-defined $\mu$-a.e.~on $A$. Then we define the \textbf{first return map to $A$} by
\begin{displaymath}
  f_{A}(x):=f^{\rho_{A}(x)}(x),
\end{displaymath}
which is well define $\mu$-a.e.~on $A$. In order to see that $f_{A}$ is a well defined dynamical system, is enough to observe that if we define
\begin{displaymath}
  A_{f}:=A\cap \left(\bigcap_{m\geq 0} \bigcup_{n\geq m} f^{-n}(A)\right),
\end{displaymath}
then by Corollary~\ref{cor:return-infinitely-many-times-Poincare} one have that $\mu(A_{f})=\mu(A)>0$, with $A_{f}\subset A$, and one can easily check that $\rho_{A}(x)<\infty$  and $f_{A}(x)\in A_{f}$, for every $x\in A_{f}$. So, this means that $f_{A}\colon A_{f}\carr$ is a well defined dynamical system. Moreover, if define the probability measure $\mu_{A}$ on the measure space $(A,\mc{B}_{A})$ by
\begin{displaymath}
  \mu_{A}(E):=\frac{\mu(E)}{\mu(A)}, \quad\forall E\in\B_{A},
\end{displaymath}
and where $\B_{A}:=\{B\in\B : B\subset A_{f}\}$, we have the following

\begin{theorem}\label{thm:inductive-map-is-mpm}
  The map $f_{A}\colon A_{f}\carr$ is a $\B_{A}$-measurable and the measure $\mu_{A}$ is $f_{A}$-invariant.
\end{theorem}

\begin{proof}
  For the sake of simplicity of the notation, from now on we shall assume that $A=A_{f}$, and for each $n\in\N$, let us consider the set $A_{n}:=\rho^{-1}(n)\subset A$.
  Observe that
  \begin{displaymath}
    A=\bigsqcup_{n\in\N} A_{n},
  \end{displaymath}
  and for each $E\subset A$, it holds
  \begin{equation}
    \label{eq:f-A-inv-on-rho-level-sets}
    f_{A}^{-1}(E)=\bigsqcup_{n\in\N} f_{A}^{-1}(E)\cap A_{n} = \bigsqcup_{n\in\N} f^{-n}(E)\cap A_{n}.
  \end{equation}

  On the other hand, one can easily check that
  \begin{equation}
    \label{eq:rho-return-times-level-sets}
    A_{n}=\rho^{-1}(n)=A\cap \bigg(\bigcap_{i=1}^{n-1} f^{-i}(X\setminus A)\bigg) \cap f^{-n}(A), \quad\forall n\in\N.
  \end{equation}

  From~\eqref{eq:rho-return-times-level-sets} we conclude $\rho$ is a measurable function and then, by~\eqref{eq:f-A-inv-on-rho-level-sets} we get that $f_{A}$ is a $\B_{A}$-measurable map.

  Finally, in order to show that $\mu_{A}$ is $f_{A}$-invariant, it is enough to show that $\mu(E)=\mu(f^{-1}_{A}(E))$, for any $E\in\B_{A}$. In order to prove that, let us fix a measurable set $E\in\B_{A}$ and observe that, by~\eqref{eq:f-A-inv-on-rho-level-sets} we get
  \begin{displaymath}
    \mu\big(f_{A}^{-1}(E)\big) = \sum_{n=1}^{\infty}\mu\big(f^{-n}(E)\cap A_{n}\big).
  \end{displaymath}
  Let us prove that the right hand side of this equality is equal to $\mu(E)$. To do that, let us start noticing that
  \begin{equation}
    \label{eq:mu-E-equal-m-f-1-E-partition-A}
    \begin{split}
      \mu(E)=\mu\big(f^{-1}(E)\big) &= \mu\big(f^{-1}(E)\cap A\big) + \mu\big(f^{-1}(E)\setminus A\big) \\
      &= \mu\big(f^{-1}(E)\cap A_{1}\big) + \mu\big(f^{-1}(E)\setminus A\big),
    \end{split}
  \end{equation}
  where the first equality follows from the fact $f^{-1}(E)=\big(f^{-1}(E)\cap A\big)\sqcup \big(f^{-1}(E)\setminus A\big)$, and the second one from $f^{-1}(E)\cap A=f^{-1}(E)\cap A_{1}$.

  Then, focusing on the last term of~\eqref{eq:mu-E-equal-m-f-1-E-partition-A} we get
  \begin{equation}
    \label{eq:mu-E-f-1-E-setminus-A-second-argument}
    \begin{split}
      \mu\big(f^{-1}(E)\setminus A\big) &= \mu\big(f^{-2}(E)\setminus f^{-1}(A)\big) \\
                                        &= \mu\Big(\big(f^{-2}(E)\setminus f^{-1}(A)\big)\cap A \Big) + \mu\Big(\big(f^{-2}(E)\setminus f^{-1}(A))\setminus A\Big)\\
      &= \mu\big(f^{-2}(E)\cap A_{2}) + \mu\Big(f^{-2}(E)\setminus\big(A\cup f^{-1}(A)\big)\Big),
    \end{split}
  \end{equation}
  where the first equality follows from the invariance of $\mu$ by $f$, the second one by the partition argument we used above to prove~\eqref{eq:mu-E-equal-m-f-1-E-partition-A}, and the last equality follows from the fact that $\big(f^{-2}(E)\setminus f^{-1}(A)\big)\cap A = f^{-1}(E)\cap A_{2}$.

  Combining equations~\eqref{eq:mu-E-equal-m-f-1-E-partition-A} and~\eqref{eq:mu-E-f-1-E-setminus-A-second-argument} and repeating inductively the last argument, one can show that
  \begin{displaymath}
    \mu(E)= \sum_{i=1}^{n} \mu\big(f^{-i}(E)\cap A_{i}\big) + \mu\Big(f^{-n}(E)\setminus\big(A\cup f^{-1}(A)\cup\cdots\cup f^{-(n-1)}(A)\big)\Big),
  \end{displaymath}
  for every $n\geq 1$.

  So, in order to prove that $\mu(E)=\mu\big(f^{-1}(E)\big)$ it is enough to show that
  \begin{displaymath}
    \mu\bigg(f^{-n}(E) \setminus \Big(\bigcup_{i=1}^{n-1} f^{-i}(A) \Big)\bigg)\leq  \mu\bigg(f^{-n}(A) \setminus \Big(\bigcup_{i=1}^{n-1} f^{-i}(A) \Big)\bigg) \to 0
  \end{displaymath}
  as $n\to +\infty$. But this easily follows from the fact that $f^{-1}(A)\setminus A, f^{-2}(A)\setminus (A\cup f^{-1}(A))\ldots$ are pairwise disjoint and hence,
  \begin{displaymath}
    1=\mu(X) \geq \sum_{n=1}^{\infty}  \mu\bigg(f^{-n}(A) \setminus \Big(\bigcup_{i=1}^{n-1} f^{-i}(A) \Big)\bigg).
  \end{displaymath}
\end{proof}


\subsection{Kakutani-Rokhlin towers}\label{sec:Kakutani-Rokhlin-towers}

In this section we shall try to revert the previous construction we performed in \S\ref{sec:induct-or-renormalization}. To do that, for the time being let $(X,\B,\mu)$ be an arbitrary measure space, $\rho\colon X\to\N_{0}$ be an integrable function and $f\colon X\carr$ be a measure preserving map.

We first onsider the new measure space $(X_{\rho},\B_{\rho},\mu_{\rho})$ where each object is defined as follows: we first define the space $X_{\rho}$ by
\begin{displaymath}
  X_{\rho}:=\bigsqcup_{n\in\N_{0}} \bigsqcup_{0\leq i\leq n} \rho^{-1}(n)\times\{i\} \subset X\times\N_{0};
\end{displaymath}
the $\sigma$-algebra $\B_{\rho}$ is defined as the one generated by the family
\begin{displaymath}
  \B_{\rho}:=\sigma\left\{A\times\{i\}\subset X\times\N_{0} : A\in\B,\ \exists n\in\N_{0},\ 0\leq i\leq n,\ A\subset\rho^{-1}(n)\right\};
\end{displaymath}
and the measure $\mu_{\rho}$ is characterized as the only measure on the measurable space $(X_{\rho},\B_{\rho})$ such that
\begin{displaymath}
  \mu_{\rho}(A\times\{i\}):={\bigg(\int_{X}\rho+1\dd\mu\bigg)}^{-1}\mu(A),
\end{displaymath}
for every generating set $A\times\{i\}$ of $\B_{\rho}$.

\begin{exercise}
  Show that $(X_{\rho},\B_{\rho},\mu_{\rho})$ is a well-defined probability space.
\end{exercise}

Then we define the measure preserving map $f_{\rho}\colon X_{\rho}\carr$ by
\begin{displaymath}
  f_{\rho}(x,i):=
  \begin{cases}
    (x,i+1), & \text{for } x\in X,\ 0\leq i<\rho(x); \\
    \big(f(x),0\big), & \text{for } x\in X,\ i=\rho(x).
  \end{cases}
\end{displaymath}


\begin{exercise}
  Show that the map $f_{\rho}\colon (X_{\rho},\B_{\rho},\mu_{\rho})\carr$ is indeed a measure preserving map. Now consider the first return map (also called induction as in \S\ref{sec:induct-or-renormalization}) of the map $f_{\rho}$ to the set $A:=X\times\{0\}$. Show that $\mu_{\rho}(A)>0$ and the first return map is ``isomorphic'' to the original $f\colon (X,\B,\mu)\carr$. Can you give a definition of \emph{isomorphism} between measure preserving maps?
\end{exercise}

\subsection{Some important exercises from~\cite[\S1.4.4]{VianaOliveiraFoundErgTh}}
\label{sec:some-import-exerc-1.4.4}

1.4.3, 1.4.4.
